考虑求 $F(x)G(x) \equiv 1 \pmod{x^n}$ 中的 $G(x)$。

若考虑分治，设我们已经知道了 $G(x)$ 的低 $x^{\frac{n}{2}}$ 项系数多项式 $G_*(x)$，看看能不能求出 $G(x)$
\[
G(x) = F^{-1}(x) \pmod{x^n}
\]
\[
G_*(x) = F^{-1}(x) \pmod{x^{\frac{n}{2}}}
\]
由于 $G(x) \equiv G_*(x) \pmod{x^{\frac{n}{2}}}$，考虑作差有
\[
G(x) - G_*(x) \equiv 0 \pmod{x^{\frac{n}{2}}}
\]
为了实现扩域，两边同时平方有
\[
\bigl(G(x) - G_*(x)\bigr)^2 \equiv 0 \pmod{x^n}
\]
也就是说
\[
G(x)^2 - 2G(x)G_*(x) + G_*(x)^2 \equiv 0 \pmod{x^n}
\]
两边乘 $F(x)$ 有
\[
G(x) \equiv 2G_*(x) - G_*(x)^2F(x) \pmod{x^n}
\]
由于 $G_1(x) = F_1^{-1}(x)$，那么之后的项可以通过递推得到。

我们发现 $G_*(x)^2F(x)$ 是不超过 $2n$ 次的多项式，但是我们计算一定会用到 $2n$ 长度的卷积，那有没有什么办法卡掉这个常数呢？

当然是有的。考虑 $G(x)$ 与 $F(x)$ 乘积的意义：
\[
G_*(x)F(x) \equiv 1 + Q(x) \pmod{x^n}
\]
其中 $Q(x)$ 为 $x^{\frac{n}{2}} \sim x^{n-1}$ 的项的组合，那么
\[
G_*(x)^2F(x) \equiv G_*(x) + Q(x)G_*(x) \pmod{x^n}
\]
其中 $G_*(x)$ 对 $\frac{n}{2}$ 以后的部分无贡献，也就是说我们实际上只需要 $Q(x)$ 与 $G_*(x)$ 卷积后的后 $\frac{n}{2}$ 项即可，前面的项对我们无意义。利用这个性质可以通过下面的操作来实现卡常：

由于 NTT 是定义在 $n = 2^k$ 序列上的循环卷积，用 $F(x)$ 与 $G_*(x)$ 线性卷积后，得到的是 $0 \sim x^{\frac{3n}{2}-1}$ 的卷积序列。进行 $n$ 长度的循环卷积时，属于 $x^{n} \sim x^{\frac{3n}{2}-1}$ 部分的贡献会计算到 $x^{0} \sim x^{\frac{n}{2}-1}$ 上。因为 $1$ 和卷积溢出部分的贡献我们都不需要，所以可以直接把前置位清空进行下一次卷积，这样卷积出来的结果取后 $\frac{n}{2}$ 就是正确的。